\section{Changing the sampling law without revealing the path}
\label{sec:sampling}

The decoding argument will use uniform matrices at two different levels.
Here we prove the distributional comparisons that permit those changes.
There are two distinct conclusions: a total-variation comparison of the
\emph{displayed functions}, and a posterior bound for a child selected after
inspecting those functions.  Neither conclusion permits disclosure of the
hidden path or of the projection mask.  All statements below hold conditional
on the complete original question tuple $E$; consequently they also hold after
averaging over $E$, uniformly in the size of the source instance.
We first replace the recursive draws at a cut by uniform draws, then compare
that law with one in which a sparse set of children is projected.  Each
comparison retains only the observation specified in its lemma.

For probability measures on a finite set, use
\[
 \TV(P,Q)=\frac12\sum_x|P(x)-Q(x)|,
 \qquad
 \chi^2(Q\Vert P)=\E_P[(dQ/dP-1)^2]
\]
when $Q\ll P$.  Cauchy--Schwarz gives
$\TV(P,Q)\le\tfrac12\sqrt{\chi^2(Q\Vert P)}$.
An observation formed by a deterministic map cannot increase total variation:
its discrepancy on an event is the discrepancy on that event's inverse image.
The same assertion holds on adjoining independent randomness and applying a
randomized map, by conditioning on that randomness.  These elementary facts
will let us retain complete assembled queries without treating their
components as independent.

\subsection{Stopping the ancestor generators}

A scalar call means one invocation of the law $\nu_D$ defined in
Section~\ref{sec:construction}, with its own fresh random choices.  Write
$b_p=2^{\ell_p}-1$ and $c_p=2R_p$.  For a fixed path prefix ending at a node
$D_p$ of level $p$, a call at a higher path node branches into exactly $c_h$
independent calls at its path child when it crosses level $h$.  Its ordinary
child contributions are uniform functions supported outside that child and
require no calls farther down the distinguished path.

\begin{lemma}[Stopped-call representation]\label{lem:stopped-calls}
Fix $E$ and the path from the root through $D_p$, including $D_p$ but no
choice of its child.  The original sampling can be generated as follows.
First generate an exterior record $\Xi_p$ consisting of the random choices
outside $D_p$'s subtree, with every recursive call reaching $D_p$ left
unevaluated.  Then generate the pending scalar calls at $D_p$, its own
buckets, and the arrays strictly below $D_p$.  Every displayed array is a
deterministic function of this record and these remaining values.

The pending ancestor calls together with the own buckets number at most
\begin{equation}\label{eq:stopped-count}
 T_p=\sum_{k=p}^d b_k\prod_{h=p+1}^k c_h.
\end{equation}
An empty product is one.  If $e_k$ additional fresh scalar draws are requested
at each level $k\ge p$, a sufficient bound is
\begin{equation}\label{eq:stopped-extra-count}
 T_p(e)=\sum_{k=p}^d(b_k+e_k)\prod_{h=p+1}^k c_h.
\end{equation}
In particular, replacing one level-$k$ bucket uses one extra draw and adds
$\prod_{h=p+1}^k c_h$ calls at the cut.  The count is independent of $n_p$ and
of all higher branching numbers.

Conditional on the full path, distinct stopped calls have independent
randomness.  Pending ancestor calls are independent of the tapes of the
actual arrays at or below the cut; own buckets are independent of the tapes
of actual descendant arrays.  Calls for the own buckets and calls pending
from ancestors are distinct even
when they sample the same function space.  The exterior record can be exposed
before sampling the continuation of the path.  Already assembled ancestor
arrays are deterministic outputs of the construction and are not independent
exterior information.
\end{lemma}

\begin{proof}
Expand each of the $b_k$ scalar buckets used for the actual matrix at $D_k$,
for every $k\ge p$.  A call at level $h>p$ is a sum of ordinary-child uniform
functions and $R_h$ products of two fresh calls at the path child.  Stop each
of those factors the first time it reaches $D_p$.  Thus a bucket originating
at level $k$ produces exactly $\prod_{h=p+1}^k c_h$ stopped calls.  Adding over
$k$ proves \eqref{eq:stopped-count}; expanding additional calls proves
\eqref{eq:stopped-extra-count}.  Ordinary children at every step are disjoint
from $D_p$'s subtree, so their contributions can be generated before any path
choice below $D_p$.

Give each recursive invocation a separate random tape, indexed by its
originating matrix, bucket and sequence of factor indices.  These indices
are distinct even if two invocations visit the same node.  The sampler's
fresh-call convention says precisely that these tapes and the tapes of
actual descendant matrices are independent, conditional on the shared path.
All unevaluated expressions are finite expressions involving sums and
pointwise products.  Filling their stopped leaves therefore reconstructs
every ancestor array deterministically.  The expression may involve many
exterior functions when higher branching is large, but its number of stopped
leaves is still \eqref{eq:stopped-count}.  This proves all the assertions.
\end{proof}

We will also use an equivalent representation after uniformizing the cut.
Instead of generating $M_{D_p}$ from $b_p$ uniform buckets, generate its
$\ell_p$ rows as separate uniform draws.  The relevant number of calls is then
\begin{equation}\label{eq:direct-cut-count}
 \widehat T_p=\ell_p+
     \sum_{k=p+1}^d b_k\prod_{h=p+1}^k c_h.
\end{equation}
We reserve $\widehat T_p$ for the displayed original call family.  If a
comparison adds $e$ independent uniform rows, where $e$ is fixed, its
enlarged count is $\widehat T_p+e$.  For instance, forming a uniform matrix
together with one rank-one step from it requires its rows and one additional
independent scalar draw, giving the count $\widehat T_p+1$.
These rows have separate tapes from the ancestor calls.  If visibility of
upper buckets will later be chosen by a row map $A$, predetermine calls for
\emph{all} buckets before choosing which are visible.  Selecting or discarding
calls afterward adds no dependence to the count.

\subsection{Finite child blocks and a hidden special child}

Fix a level-$p$ cut $D=D_p$ and an integer $t_{\mathrm{cut}}$ bounding all scalar draws
retained at that cut.  For a child $C$ of $D$, put $V_C=H_C(E)^2$ and define
the following finite product space:
\begin{equation}\label{eq:raw-child-block}
 \mathcal X_C=V_C^{t_{\mathrm{cut}}}\times
  \prod_{\substack{T\subseteq C\\T\text{ nonleaf}}}
       (K_{\operatorname{level}(T)}\otimes H_T(E)).
\end{equation}
Here $T\subseteq C$ means that $T$ is $C$ or a descendant of $C$.
A block $X_C$ records a separate child contribution to each of the $t_{\mathrm{cut}}$ cut
draws and the actual arrays throughout this child's subtree.  The factors
used internally to generate a special contribution are not included in
$X_C$.  Each recorded contribution is just a function in $V_C$.

Let $P_C$ be uniform on $\mathcal X_C$.  To construct independent uniform cut
draws, sample these blocks independently and, for each draw, add its child
contributions.  This uses a latent product of spaces even though the actual
spaces $V_C$ share the constant functions.  Indeed the map
\[
 \prod_{C\text{ child of }D}V_C\longrightarrow H_D(E),
 \qquad (u_C)_C\longmapsto\sum_Cu_C
\]
is a surjective linear map.  Every fiber is a coset of its kernel, so the
image of uniform measure is uniform.  Separate copies give independent
uniform cut draws, independently of all descendant arrays.  No direct-sum
decomposition, or recovery of the summands from their sum, is asserted.

There is a uniform finite density bound for any law on $\mathcal X_C$.
For clarity, if the source's legal assignment domains have size at most $a$,
put $N_0=1$ and $N_k=\prod_{h=1}^k n_h$.  Every space on a level-$k$ subtree
has dimension at most $a^{N_k}$.  Thus, for children at level $p-1$,
\begin{equation}\label{eq:block-dimension-bound}
 \dim\mathcal X_C\le
 D_p(t_{\mathrm{cut}}):=a^{N_{p-1}}\left(t_{\mathrm{cut}}+
  \sum_{k=1}^{p-1}\ell_k\prod_{h=k+1}^{p-1}n_h\right).
\end{equation}
In particular we may use $L_p(t_{\mathrm{cut}})=2^{D_p(t_{\mathrm{cut}})}$.  Every probability law $Q_C$ on
$\mathcal X_C$ has density $q_C=dQ_C/dP_C\le L_p(t_{\mathrm{cut}})$, since each $P_C$ atom has
mass $2^{-\dim\mathcal X_C}$.  This bound uses the source alphabet, the fixed row
and repetition counts through $t_{\mathrm{cut}}$, and branching below $p$.  It uses neither
$n_p$, higher branching, source question identities, nor the source instance
size.  Any larger fixed common bound will also be denoted by $L_p$.

\begin{lemma}[One hidden child]\label{lem:cut-tv}
In the original experiment, expose $E$, the path through $D_p$, and the
exterior record of Lemma~\ref{lem:stopped-calls}.  Forget the entire
continuation of the path below $D_p$.  Retain any fixed family of stopped
scalar values, the own buckets at $D_p$, and all actual arrays below $D_p$.
Let $\mathsf U_p$ be the \emph{unrestricted experiment}: replace the retained
cut scalar values by independent uniform draws in $H_{D_p}(E)$, and replace
all descendant arrays by independent uniform arrays in their respective
spaces, retaining the same exterior record and reconstruction expressions.
Then
\begin{equation}\label{eq:cut-tv-bound}
 \TV(\mathsf A_p,\mathsf U_p)\le
 \varepsilon_p^{\mathrm{cut}}:=\frac{L_p}{2\sqrt{n_p}}
\end{equation}
for the retained values and for every deterministic observation of them.
Here $\mathsf A_p$ denotes the corresponding retained original law, and
$L_p$ can be chosen using \eqref{eq:block-dimension-bound} with the count in
\eqref{eq:stopped-extra-count}.

One may also append a fixed number of independent uniform cut rows, or a
fixed number of fresh bucket replacements, with the appropriate increased
count.  The bound includes the reconstructed ancestor arrays and candidate
edge tests at levels at least $p$, with independent quotient directions.
It does not include a tag identifying the next path child, any later path
indices, or the internal factor tapes generating a special child block.
\end{lemma}

\begin{proof}
Condition on $E$, the prefix and the exterior record.  Let $J$ be the next
path child, uniform among $n=n_p$ children.  Given $J=C$, every block other
than $X_C$ has law $P_{C'}$: all its scalar contributions are ordinary-child
uniforms, and all its subtree matrices are off the path and uniform.  They
are independent of one another and of $X_C$.  Within $X_C$, the scalar
contributions are sums of $R_p$ products from fresh factor calls, and its
subtree arrays use the continued path.  Average over that continuation to
obtain a law $Q_C$ on \eqref{eq:raw-child-block}.  Correlations among the
coordinates of this one block cause no difficulty.  Its density $q_C$ is
bounded by $L_p$.

Relative to $P=\prod_CP_C$, the original raw-block law has density
\[
 Z=\frac1n\sum_Cq_C(X_C).
\]
Under $P$ these summands are independent and have expectation one.  Therefore
\[
 \E_P[(Z-1)^2]
   =\frac1{n^2}\sum_C\E_{P_C}[(q_C-1)^2]
   \le\frac{L_p^2}{n}.
\]
The total-variation bound follows by Cauchy--Schwarz.  Summing the child
contributions and reconstructing all pending expressions is the same
deterministic map in both laws, so it preserves the bound.  It remains valid
conditional on every exterior record and hence on averaging that record.
Fresh bucket replacements merely enlarge the finite family of scalar calls;
independent uniform cut rows add uniform coordinates to every block.  Both
are covered by the same argument.  A revealed value of $J$ would replace the
average density by $q_J$ and would lose the factor $n^{-1/2}$, which explains
the required information restriction.
\end{proof}

In $\mathsf U_p$, the own buckets are separate from the pending ancestor
calls and the descendant matrices.  The matrix
$M=\sum_{v\in K_p\setminus\{0\}}v h_v$ is uniform: the linear map from its
uniform buckets onto $K_p\otimes H_{D_p}(E)$ is surjective, since the standard
basis vectors occur among the indices.  Replacing $h_a$ by an independent
uniform $h_a^*$, for an independent nonzero $a\in K_p$, gives
\begin{equation}\label{eq:unrestricted-step}
 M^*=M+a h,\qquad h=h_a+h_a^*.
\end{equation}
Conditional even on all old buckets, $h$ is uniform and independent of the
old data.  Thus this is exactly the uniform rank-one step, and $M$ is
independent of the other displayed arrays and the pending ancestor calls.
This conclusion would not follow by reusing a bucket as an ancestor call.

\subsection{Sparse projection and its posterior}

The next comparison is a form of sub-code covering: a sparse restriction
leaves the lifted query law nearly unchanged.  This method appears in
Khot and Safra~\cite[Section~3.4]{KhotSafra2013}, with advice-conditioned
thinning estimates in Dinur, Khot, Kindler, Minzer and Safra
\cite[Lemma~5.5, full version]{DinurKhotKindlerMinzerSafra2018} and
Fei, Minzer and Wang~\cite[Propositions~4.7--4.8 and~5.5]{FeiMinzerWang2026}.
Here the independent units are entire child subtrees, whose records include
all stopped-call contributions and displayed descendant arrays.  We prove
the comparison and posterior estimate for these finite blocks.

At a level-$i$ cut $D_i$, independently mark each child with probability
$\beta_i=n_i^{-2/3}$; take $n_i>1$ to be a cube.  At every leaf in a marked
child, sample a source projection conditional on its left question, and
replace that question by its right question.  Leave all other questions
unchanged, and call the resulting tuple $O$.  Pullback along these
projections identifies $H_T(O)$ with a subspace of $H_T(E)$ and preserves
pointwise multiplication.  Define $\mathsf M_i$, the \emph{modified
experiment}, by stopping all ancestor calls at $D_i$ and making their values
independent uniform draws in $H_{D_i}(O)$.  Also sample $M_{D_i}$ and all its
descendant matrices independently uniformly in the corresponding projected
spaces.  Reconstruct all ancestors by the original stopped expressions.
Path choices below $i$ are unused.  Independent extra cut rows and bucket
replacements are treated in the same way.  The original unmodified law will
be denoted by $\mathsf A$.

\begin{lemma}[Sparse projection comparison]\label{lem:sparse-tv}
Fix $E$, the path through $D_i$ and an exterior record generated without the
projections.  Compare $\mathsf M_i$ and $\mathsf U_i$ as laws of lifted
functions on $E$, forgetting the marks, the projection choices, and the
projected-question tuple $O$.  All raw child summands and descendant arrays
may be retained.  For a bound $L_i$ from
\eqref{eq:block-dimension-bound},
\begin{equation}\label{eq:sparse-tv-bound}
 \TV(\mathsf M_i,\mathsf U_i)
 \le \varepsilon_i^{\mathrm{spr}}
 :=\frac12\sqrt{(1+\beta_i^2L_i^2)^{n_i}-1}.
\end{equation}
The bound holds for any fixed number of independent extra cut rows and fresh
higher-level bucket draws, with their enlarged call count.  In particular it
holds for observations containing the assembled ancestor arrays and both
candidate tests at levels $i<j$.

Consequently the modified experiment and the original experiment have
observation distance at most
$\varepsilon_i^{\mathrm{spr}}+\varepsilon_i^{\mathrm{cut}}$
when the path below $i$ is forgotten.  For the further comparison at a level
$j>i$, require the observation to be a common deterministic function of the
exterior record $\Xi_j$, the retained scalar values at $D_j$ (including its
own buckets and any extra cut draws), all displayed arrays strictly below
$D_j$, and independent auxiliary choices.  Here $E$ and the path through
$D_j$ are fixed.  Such an observation may include reconstructed ancestors
and descendant arrays with their tree locations, but it cannot use the
hidden continuation of the path or the internal decompositions of the
retained scalar values.  In particular, a selected lower-level test is
permitted only if it is determined by these data without using the hidden
path.  The distance from the corresponding unrestricted-$j$ observation is
then at most
\begin{equation}\label{eq:two-cut-triangle}
 \varepsilon_i^{\mathrm{spr}}+
 \varepsilon_i^{\mathrm{cut}}+\varepsilon_j^{\mathrm{cut}}.
\end{equation}
\end{lemma}

\begin{proof}
For a child $C$ of $D_i$, let $P_C$ be uniform on the unprojected block
\eqref{eq:raw-child-block}.  Conditional on marking $C$, average over all its
projection choices and lift the independent projected uniform contributions
and arrays to $E$; call the resulting law $Q_C$.  It is a law on the same
finite set, so $q_C=dQ_C/dP_C\le L_i$.  Different children use independent
marks, projections and tapes.  Thus their lifted blocks have product law
\[
 \prod_C\bigl((1-\beta_i)P_C+\beta_iQ_C\bigr),
\]
whose density with respect to $\prod_CP_C$ is
$\prod_C(1+\beta_i(q_C-1))$.  Squaring and using independence yields the
exact identity
\[
 1+\chi^2\!\left(
   \prod_C((1-\beta_i)P_C+\beta_iQ_C)
      \,\middle\Vert\,\prod_CP_C\right)
  =\prod_C\bigl(1+\beta_i^2\chi^2(Q_C\Vert P_C)\bigr)
  \le(1+\beta_i^2L_i^2)^{n_i}.
\]
This proves \eqref{eq:sparse-tv-bound}.  Uniform cut rows are obtained by the
same addition map in both laws.  A uniform matrix can be represented either
by independent rows or by independent uniform buckets, as proved above;
choose the latter representation when retaining a bucket replacement.
Additional calls simply add coordinates to each finite block.

All higher arrays, and their quotients in a candidate test, are functions of
the blocks, exterior record and independent test directions.  Data processing
therefore proves their bounds, without conditioning on their realized values.
The comparison with $\mathsf A$ is the triangle through $\mathsf U_i$ and
Lemma~\ref{lem:cut-tv}.  For \eqref{eq:two-cut-triangle}, take the common
observation just specified.  First average the preceding
$i$-cut comparison over the prefix between levels $j$ and $i$, and then apply
Lemma~\ref{lem:cut-tv} at $j$ to the original law.  A pending scalar value at
$j$ is an output of the $i$-cut reconstruction and is a retained cut draw in
the $j$-cut comparison, so it can be included throughout.  The internal
factor choices within its subtree cannot be included in the latter
comparison.  This establishes precisely the stated triangle.
\end{proof}

\paragraph{Common observations in the three comparisons.}\label{sec:law-guide}
The laws just defined serve different purposes, and each comparison forgets
different information:
\begin{description}[style=nextline,leftmargin=1.5em,font=\normalfont\bfseries]
\item[Original to unrestricted at one cut.]
Lemma~\ref{lem:cut-tv} replaces the cut draws and descendant arrays by
uniforms after forgetting the path below that cut.  The observation may
include the ancestors reconstructed from those draws.
\item[Modified to unrestricted at the lower cut.]
Lemma~\ref{lem:sparse-tv} compares lifted functions on $E$ after forgetting
$O$, the projection choices, and the marks.  Raw child contributions and
all displayed descendant arrays may be retained.
\item[Modified lower cut to unrestricted upper cut.]
For $j>i$, Equation~\eqref{eq:two-cut-triangle} applies to a common function
of the upper-cut record specified in Lemma~\ref{lem:sparse-tv}.  That record
retains pending scalar values there, but not their internal decompositions,
and all displayed descendant arrays.  The observation cannot select among
those arrays using the hidden path.
\end{description}

The sparse product calculation also gives a posterior bound for a child
selected after seeing the lifted matrices.  This is stronger than a bound
for a child fixed in advance, and will allow the lower decoder to choose a
witness from its observed data.

\begin{lemma}[Posterior bound for adaptive witnesses]\label{lem:posterior}
In $\mathsf M_i$, condition on $E$, the path through $D_i$ and a permissible
exterior record.  Let $X=(X_C)_C$ consist of all lifted raw child blocks,
including any fixed number of independent uniform cut rows.  Let
$\mathcal G$ be the sigma-algebra generated by these data and additional
independent random choices, such as a test direction or fixed full strategy
tapes.  It contains neither projected-question metadata nor projection marks.
For every child $C$, almost surely,
\begin{equation}\label{eq:posterior-bound}
 \Prob(C\text{ marked}\mid\mathcal G)
  \le q_i:=\frac{\beta_i L_i}{1-\beta_i}.
\end{equation}
If a set $\mathcal C(X)$ of at most $m$ children is selected measurably from
these data, then
\begin{equation}\label{eq:adaptive-child-bound}
 \Prob(\mathcal C(X)\text{ contains a marked child})\le m q_i.
\end{equation}
Both assertions remain valid for a coarsening of this information.

In particular, let $F_0,F_1$ be bilinear forms on $H_{D_i}(E)$ determined by
these lifted data.  Then
\begin{equation}\label{eq:adaptive-bilinear-error}
 \Prob\bigl(F_0\ne F_1,\ 
   F_0|_{H_{D_i}(O)\times H_{D_i}(O)}
       =F_1|_{H_{D_i}(O)\times H_{D_i}(O)}\bigr)
 \le 2q_i,
\end{equation}
where the restrictions in this display mean restriction of both arguments
to $H_{D_i}(O)$.
No part of the statement conditions on strategy success, a right output,
a slice event, or a witness-matching event.
\end{lemma}

\begin{proof}
Conditional on the fixed exterior data, child independence and Bayes's
formula give, simultaneously at every lifted block tuple with positive
probability,
\[
 \Prob(C\text{ marked}\mid X)
 =\frac{\beta_iq_C(X_C)}{1-\beta_i+\beta_iq_C(X_C)}
 \le\frac{\beta_iL_i}{1-\beta_i}.
\]
Independent extra choices do not change this conditional formula.  For an
adaptively selected set, condition on $\mathcal G$ first and use
\[
 \E\left[\sum_{C\in\mathcal C(X)}
             \one_{\{C\text{ marked}\}}\ \middle|\ \mathcal G\right]
 \le |\mathcal C(X)|q_i\le m q_i.
\]
Taking expectations and a union bound proves
\eqref{eq:adaptive-child-bound}; the tower property proves the coarsened
versions.  There is no factor for the total number of available children.

For the bilinear consequence, write $\Delta=F_0-F_1$.  Since
$H_{D_i}(E)=\sum_C H_C(E)^2$, if $\Delta\ne0$, then there are children
$C_1,C_2$ and functions $x\in H_{C_1}(E)^2$, $y\in H_{C_2}(E)^2$ with
$\Delta(x,y)=1$.  Indeed decompose a pair on which $\Delta$ is nonzero into
such sums and expand bilinearly; at least one summand is nonzero.  Fix orders
on the finite spaces and choose the first such witness from the lifted data.
The children may coincide.  If neither is marked, their spaces are unchanged
and lie in $H_{D_i}(O)$, so the two restricted forms cannot be equal.  The
bad event in \eqref{eq:adaptive-bilinear-error} is therefore contained in the
event that one of at most two adaptively selected children is marked.
Apply \eqref{eq:adaptive-child-bound} with $m=2$.
\end{proof}

The distinction from success conditioning is substantive: after the
unconditional bad-event probability is bounded, it may be subtracted from
any success probability.  The posterior estimate itself need not hold within
that success event.  Exposing a direction $a_i$ to the analyst does not
provide it to the left strategy.  Likewise, adding assembled ancestor values to the
lifted blocks changes no information, because they are functions of those
blocks and $\Xi_i$; treating those values as independently sampled exterior
data would not justify the product law used above.

\subsection{Bias that is independent of dimension}

Total variation above is obtained by increasing branching.  The upper
Fourier argument instead requires a small-bias bound before any slice
conditioning, obtained by increasing the product counts.  We use the
character-wise notion of small bias of Naor and Naor~\cite{NaorNaor1993};
the required sampler estimate is proved below.  For a law $\mu$
on a finite binary vector space $V$, call it \emph{quarter-balanced} if every
nonzero $f\in V^*$ satisfies
$\Prob_{x\sim\mu}(f(x)=b)\ge1/4$ for both $b\in\F$.

\begin{lemma}[Scalar bias and independent hidden buckets]\label{lem:bias}
There is an absolute integer $R_{\min}$ such that the following holds if
$R_p\ge R_{\min}$ at every nonleaf level.  Fix the domains and the full path
in the original sampler, or fix $E,O$, all projection choices, and the path
through a modified cut in $\mathsf M_i$.  Every scalar law is quarter-balanced.
For each nonzero $f\in H_{D_p}^*$ at an original path node, or
$f\in H_{D_p}(O)^*$ at a node strictly above the modified cut,
\begin{equation}\label{eq:scalar-bias}
 \left|\E_{h\sim\nu_{D_p}}(-1)^{f(h)}\right|
       \le (7/8)^{R_p}.
\end{equation}
At leaves, and at or below the uniformized cut, uniform scalar draws have
zero bias on nonzero functionals.

More generally, fix $j>i$, a linear map $A:K_j\to\F^r$, and $W=\ker A$.
In the modified experiment let $h_v$, $v\ne0$, be the separate level-$j$
buckets and set
\[
 X=\sum_{v\in W\setminus\{0\}}v h_v
     \in W\otimes H_{D_j}(O).
\]
Conditional on the domain and path data just specified, all other actual
matrices, and all buckets with $Av\ne0$, every nonzero character of $X$ has
bias at most $(7/8)^{R_j}$.  This remains so on revealing original questions
and any other independent sampling tapes.  Its law is determined by $O$,
the path through the cut, $A$, and the fixed sampling parameters; it does
not require the original questions or the missing part of a lower matrix.
Revealing a full lower matrix here is an analytical exposure; the right
decoder still receives only its prescribed quotient.  These conclusions
are not asserted after conditioning on a column slice, on successful
decoding, or on cleanliness.
\end{lemma}

\begin{proof}
First suppose $g,h$ are independent draws from quarter-balanced laws on a
binary space $V$, and $B:V\times V\to\F$ is a nonzero bilinear form.  The
linear map $g\mapsto B(g,\cdot)$ has proper kernel.  Some nonzero linear
functional on $V$ vanishes on that kernel; quarter balance therefore implies
\[
 \Prob(B(g,\cdot)\ne0)\ge1/4.
\]
For each such $g$, quarter balance of the independent $h$ implies
$\Prob(B(g,h)=b\mid g)\ge1/4$ for each bit $b$.  Hence each bit of $B(g,h)$
has probability at least $1/16$, and
$|\E(-1)^{B(g,h)}|\le7/8$.

Now apply induction from the leaves, or upward from the modified uniform
cut.  Let $f$ be nonzero on the space at a path node $D_p$.  If its
restriction to the square space of some ordinary child is nonzero, the
corresponding uniform summand makes $f(h)$ an exactly uniform bit.  Otherwise
$f$ vanishes on every ordinary child's square space.  Since these spaces
together with the special child's square space span $H_{D_p}$, its
restriction to the latter is nonzero.  The bilinear form
\[
 B(g,h)=f(gh)
\]
on that child's row space is then nonzero, since pointwise products span its
square space.  The induction hypothesis makes the two independent factor
laws quarter-balanced.  Each product thus has sign expectation of absolute
value at most $7/8$, and the $R_p$ independent products give
\eqref{eq:scalar-bias}.  Choose $R_{\min}$ so that
$(7/8)^{R_{\min}}\le1/2$.  Then the resulting scalar law is again
quarter-balanced, closing the induction.  Uniform cut and leaf laws provide
the initial step.  The argument uses only whether linear or bilinear maps
are nonzero, and no dimension enters its bound.

For the hidden sum, identify a character with a functional
$\Phi\in(W\otimes H_{D_j}(O))^*$.  If $\Phi\ne0$, then for some nonzero
$v\in W$ the functional $h\mapsto\Phi(v\otimes h)$ is nonzero; simple
tensors span the tensor product.  Conditional on the stated information,
every hidden bucket retains its own fresh tape, independently of the other
hidden buckets.  This follows directly from the call indexing in
Lemma~\ref{lem:stopped-calls}; higher matrices use distinct calls even where
they visit the same spaces.  Consequently
\[
 \E(-1)^{\Phi(X)}=
 \prod_{v\in W\setminus\{0\}}
    \E(-1)^{\Phi(v\otimes h_v)}.
\]
One factor has absolute value at most $(7/8)^{R_j}$ and the others at most
one.  If $W=0$ there are no nonzero characters and the assertion is vacuous.
Every hidden bucket is generated on the projected domain $O$ by fresh
recursion ending with uniform draws at the cut.  That recipe uses neither
$E$ nor the values of other actual matrices, proving the asserted law and
its invariance under the permissible additional exposure.
\end{proof}

The order of accuracy choices follows from this proof.  At descendants one
needs only the absolute quarter-balance guarantee, obtained by the common
$R_{\min}$.  Once an upper level $j$ requires a tolerance $\lambda_j>0$, it
suffices to choose $R_j$ with $(7/8)^{R_j}\le\lambda_j$.  No requested upper
accuracy imposes a finer accuracy at a lower level.  After all row counts,
product counts and the source alphabet are fixed, choose branching upward.
At cut $p$, $L_p$ is already fixed by lower branching, and
\[
 \varepsilon_p^{\mathrm{cut}}\longrightarrow0,
 \qquad
 \varepsilon_p^{\mathrm{spr}}\longrightarrow0,
 \qquad
 \frac{2n_p^{-2/3}L_p}{1-n_p^{-2/3}}\longrightarrow0.
\]
For the middle limit, use
$(1+n_p^{-4/3}L_p^2)^{n_p}\le\exp(L_p^2n_p^{-1/3})$.
Enlarging higher branching changes only the exterior record and the
reconstruction map, so all three bounds already secured at $p$ persist.
